\chapter{Connection to Synapse and maubot}
\label{ch:synapse}

\section{The bot's identity}

The bot's Matrix account \cmd{@adminbot} is the only privileged identity in the system. It
has Synapse's server-admin flag; professors and TAs never get credentials of any kind.
maubot holds the account's access token and device (created once with \cmd{mbc auth}, see
the setup guide). All requests go to \cmd{http://synapse:8008} on the Docker network, so the
IP allowlist that host nginx applies to \cmd{/_synapse/admin} doesn't affect the bot.

The plugin doesn't need Synapse's registration shared secret: accounts are created with the
admin API and the bot's own token.

\section{Synapse admin API}

All calls go through mautrix (\cmd{client.api.request} with \cmd{SynapseAdminPath}):

\begin{center}
\small
\begin{tabularx}{\textwidth}{@{}>{\raggedright\arraybackslash}p{6.4cm}L@{}}
\toprule
\textbf{Endpoint} & \textbf{Used for} \\
\midrule
\cmd{GET /_synapse/admin/v2/users/<id>} & does an account exist; is the bot a server admin (startup) \\
\cmd{PUT /_synapse/admin/v2/users/<id>} & create accounts from signup links (password, display name, email threepids) and professor-bot accounts (\cmd{user_type: bot}); reactivate a deleted bot's account. The endpoint also \emph{modifies} existing users, so the bot checks first that the account doesn't exist. \\
\cmd{POST /_synapse/admin/v1/reset_password/<id>} & reset links (\cmd{logout_devices: true}) \\
\cmd{POST /_synapse/admin/v1/deactivate/<id>} & deleting a professor bot \\
\cmd{POST /_synapse/admin/v1/users/<id>/override_ratelimit} & the bot exempts itself at startup (0/0 = no limit); optionally professor bots \\
\cmd{POST /_synapse/admin/v1/join/<room>} & force-join new users into their course after signup \\
\cmd{GET /_synapse/admin/v1/users/<id>/joined_rooms} & \cmd{!bot info} (room count) \\
\bottomrule
\end{tabularx}
\end{center}

A professor bot's access token comes from a normal password login
(\cmd{POST /_matrix/client/v3/login}) with a random password that is thrown away
afterwards. The admin \enquote{login as user} API isn't used because it doesn't create a
real device, which end-to-end encryption needs.

\section{Client-server API}

Everything else is ordinary client-server API as the bot user: \cmd{createRoom} (spaces and
rooms with their initial state), state events (\cmd{m.space.child}, \cmd{m.room.power_levels},
names, topics), invites, kicks, member lists, sending messages, reactions and files, and the
\cmd{m.direct} account data.

\subsection{How rooms are built}

\begin{itemize}
  \item A course is a space (\cmd{creation_content.type = m.space}) with history visibility
        \cmd{shared}. Rooms get an \cmd{m.space.parent} pointing at the space; the space gets
        \cmd{m.space.child} entries with an \cmd{order} key so the order is stable.
  \item Open course rooms use the join rule \cmd{restricted} with an allow rule for members
        of the space, so accepting one invitation to the space opens all of them. This needs
        room version 8 or newer; the server's default (v10 on Synapse 1.157) is used.
  \item Staff rooms and group rooms are \cmd{invite}-only.
  \item Rooms are created encrypted (\cmd{m.megolm.v1.aes-sha2}) unless a template says
        otherwise.
\end{itemize}

\subsection{Power levels}

Every room the bot creates starts from the same power-level content. Templates only store
deviations from it.

\begin{center}
\small
\begin{tabular}{@{}ll@{\qquad}ll@{}}
\toprule
\multicolumn{2}{@{}l}{\textbf{Actions}} & \multicolumn{2}{l@{}}{\textbf{Users}} \\
\midrule
\cmd{users_default}, \cmd{events_default} & 0 & course owner & 100 (\cmd{power_levels.owner}) \\
\cmd{state_default}, \cmd{invite}, \cmd{kick} & 50 & TAs & 50 (\cmd{power_levels.ta}) \\
\cmd{ban}, \cmd{redact} & 50 & \cmd{@adminbot} & 100 \\
\cmd{m.room.power_levels}, \cmd{m.room.encryption} & 100 & students & 0 \\
\cmd{m.space.child}/\cmd{parent}, room tag & 100 & & \\
\bottomrule
\end{tabular}
\end{center}

Archiving raises \cmd{events_default} to the TA level, so only staff can still post.
Ownership and TA status live in the bot's database (\cmd{course_staff}), never in power
levels: the power levels follow from the database, not the other way round.

\subsection{The room tag}

Every room and space the bot creates carries a state event of type
\cmd{ch.unibas.dmi.marvin.room}:

\begin{codefile}[state event ch.unibas.dmi.marvin.room]
{"kind": "group", "course": "CS101", "group": 3, "members": "staff"}
\end{codefile}

\cmd{kind} is \cmd{space}, \cmd{room} or \cmd{group}; \cmd{members} is \cmd{all} or
\cmd{staff}. Templates, groups, archive, TA removal and force-join rely on it, and a room with
this tag is never treated as a DM.

\subsection{Direct chats}
\label{sec:dm-rule}

Matrix has no server-side flag for \enquote{direct chat}. The bot treats a room as a DM when
exactly two users are joined, the bot and the sender, nobody else is invited, and the room
doesn't carry the bot's room tag. It checks this on every command (no caching), because
someone could be invited later. The audit room never counts as a DM.

Notifications (bot approved, password reset, professor bot ready) go to the user's most
recent DM, taken from the bot's \cmd{m.direct} account data, which the router updates for
every command it accepts. Without one, the bot creates a new encrypted DM.

\section{End-to-end encryption}

The bot's maubot client has its own device and an Olm/Megolm crypto store in maubot's
database, so it can read encrypted DMs and post into encrypted rooms (\cmd{!groups post},
notifications). Uploaded files in encrypted rooms are decrypted on download, and files the
bot sends (CSV, YAML) are encrypted before upload.

\begin{warning}[Crypto pickle key]
The main maubot still uses the default \cmd{crypto_db_pickle_key} (infrastructure item I3,
deliberately deferred). Changing it later invalidates the bot's crypto sessions; the client
then has to be re-created and re-verified. The userbots maubot has its own random key.
\end{warning}

\section{Rate limits}

At startup the bot sets a rate-limit override for itself, which Synapse applies to all its
rate limiters. If that fails, the startup check \cmd{ratelimit} reports it, and jobs still
pause \cmd{jobs.delay_seconds} between items. A rate-limited request fails with \cmd{429
M_LIMIT_EXCEEDED}; mautrix doesn't retry it, so the user sees an error rather than a delay.

\section{The signup web app}

maubot serves the plugin's web app under \cmd{/_matrix/maubot/plugin/<instance>/}. The
plugin registers \cmd{GET} and \cmd{POST} for \cmd{/signup} and \cmd{/reset}. Host nginx
forwards only this prefix (\cmd{location ^~ ...}, which wins over the general \cmd{/_matrix}
regex) and sets \cmd{X-Forwarded-For}. The public link base comes from the plugin setting
\cmd{signup.public_base_url}, not from maubot's \cmd{public_url}, which points at the SSH
tunnel.

\section{The userbots maubot}

Professor bots are managed through the maubot management API of the second maubot
(\cmd{user_bots.maubot_url}, default \cmd{http://userbots:29316}), with a login whose token
is cached and renewed on 401:

\begin{center}
\small
\begin{tabularx}{\textwidth}{@{}>{\raggedright\arraybackslash}p{6.4cm}L@{}}
\toprule
\textbf{Endpoint (under \cmd{/_matrix/maubot/v1})} & \textbf{Used for} \\
\midrule
\cmd{POST /auth/login} & token for the following calls \\
\cmd{POST /plugins/upload?allow_override=true} & upload a plugin (after the ownership check) \\
\cmd{POST /client/new}, \cmd{PUT /client/<id>} & the bot account's client (token, device, autojoin) \\
\cmd{POST /client/<id>/generate_recovery_key}, \cmd{/verify} & cross-sign the bot's device \\
\cmd{PUT /instance/<id>} & create, start/stop, change plugin or config \\
\cmd{GET /clients}, \cmd{GET /instances} & startup guard, kill switch \\
\cmd{DELETE /instance/<id>}, \cmd{/client/<id>}, \cmd{/plugin/<id>} & deleting a bot \\
\bottomrule
\end{tabularx}
\end{center}

Provisioning runs as a job and is safe to repeat after an interruption: an existing account
is reset or reactivated, an existing client is updated instead of created.

After creating the client, the job cross-signs the bot's new device:
\cmd{generate_recovery_key} creates cross-signing keys and signs the device, so Element shows
the bot as verified by its owner, and clients that only share room keys with cross-signed
devices still include it. The recovery key is kept in the table \cmd{bot_recovery_key}. If
the account already has keys (a reactivated bot, an interrupted job), maubot answers
\cmd{409}, and the job signs the device with \cmd{/verify} and the stored key instead. A
failure doesn't stop provisioning; the owner gets a note and the audit room the reason.
